\section{Das Artefakt: Living Knowledge Graph}

Das zentrale Artefakt dieser Arbeit ist ein System, das eine Codebase als strukturierten Wissensgraphen repräsentiert und diesen als proaktive Kontextinfrastruktur für Coding Agents bereitstellt. \emph{Living} bezeichnet dabei die Eigenschaft, dass der Graph nach jedem Commit automatisch aktualisiert wird und damit dauerhaft den aktuellen Codestand abbildet (siehe Schicht~2). Das System gliedert sich in drei Schichten; Abbildung~\ref{fig:architecture} gibt einen Überblick.

\begin{figure}[htbp]
\centering
\begin{tikzpicture}[
  node distance=0.65cm,
  mybox/.style={
    draw, rounded corners=3pt,
    minimum width=12cm, align=center,
    inner sep=9pt, font=\small
  },
  arr/.style={-{Stealth[length=5pt,width=4pt]}, thick, gray!55},
  lbl/.style={font=\scriptsize\itshape, text=gray!65}
]

\node[mybox, fill=blue!8] (codebase) {%
  \textbf{Quellcode-Repositories}\quad
  \texttt{.ts}\;/\;\texttt{.tsx}\;/\;\texttt{.py}\;/\;\texttt{.js}};

\node[mybox, fill=orange!10, below=of codebase] (s1) {%
  \textbf{Schicht 1: Ingestion (tree-sitter)}\\[3pt]
  Knoten:\enspace\texttt{File},\enspace\texttt{Function},\enspace\texttt{Class}%
  \qquad
  Kanten:\enspace\texttt{IMPORTS},\enspace\texttt{CALLS},\enspace\texttt{EXTENDS},\enspace\texttt{DEFINED\_IN}};

\node[mybox, fill=orange!14, below=of s1] (s2) {%
  \textbf{Schicht 2: Living Graph (Neo4j)}\\[3pt]
  Git post-commit Hook\enspace$\cdot$\enspace
  inkrementelle Synchronisation\enspace$\cdot$\enspace
  Git-History-Metriken};

\node[mybox, fill=orange!19, below=of s2] (s3) {%
  \textbf{Schicht 3: Context Assembly}\\[3pt]
  Entity Extraction\;$\to$\;Seed-Lookup\;$\to$\;
  Graph-Traversal\;$\to$\;Ranking\;$\to$\;Entscheidungspfad};

\node[mybox, fill=teal!10, below=of s3] (mcp) {%
  \textbf{MCP-Server}\quad
  \texttt{get\_context(task)}\enspace
  \texttt{inspect\_graph(node\_id)}\enspace
  \texttt{update\_graph()}};

\node[mybox, fill=green!8, below=of mcp] (agent) {%
  \textbf{Coding Agent}\quad
  Claude Code\;$\cdot$\;Cursor\;$\cdot$\;weitere MCP-kompatible Clients};

\draw[arr] (codebase) -- node[lbl, right=5pt]{tree-sitter Parser} (s1);
\draw[arr] (s1)       -- node[lbl, right=5pt]{Graph aufbauen} (s2);
\draw[arr] (s2)       -- node[lbl, right=5pt]{Graphabfrage via Cypher} (s3);
\draw[arr] (s3)       -- node[lbl, right=5pt]{MCP JSON-RPC} (mcp);
\draw[arr] (mcp)      -- node[lbl, right=5pt]{Kontext-Block und Entscheidungspfad} (agent);

\end{tikzpicture}
\caption{Systemarchitektur des \emph{Living Knowledge Graph}. Der Coding Agent erhält
über den MCP-Server proaktiv assemblierten, aufgabenspezifischen Kontext, ohne die
Codebase reaktiv zu explorieren.}
\label{fig:architecture}
\end{figure}


\subsection*{Schicht 1: Ingestion und Graph-Aufbau}

Ein auf \emph{tree-sitter} basierender Parser verarbeitet Python-, JavaScript- und TypeScript-Quellcode und extrahiert daraus einen strukturierten Graphen in Neo4j:

\begin{tabularx}{\linewidth}{@{}lX@{}}
  \toprule
  \textbf{Knotentyp} & \textbf{Bedeutung} \\
  \midrule
  \texttt{File}      & Quellcode-Datei mit Pfad und Sprache \\
  \texttt{Function}  & Funktion oder Methode mit Signatur und Zeilenbereich \\
  \texttt{Class}     & Klasse oder Interface \\
  \texttt{TestCase}  & Testfunktion mit Referenz zum getesteten Symbol \\
  \texttt{Schema}    & Daten- oder API-Schema (z.\,B. Sanity- oder Pydantic-Schema) \\
  \bottomrule
\end{tabularx}

\vspace{0.5em}

\begin{tabularx}{\linewidth}{@{}lX@{}}
  \toprule
  \textbf{Kantentyp} & \textbf{Bedeutung} \\
  \midrule
  \texttt{IMPORTS}   & Datei importiert ein Symbol aus einer anderen Datei \\
  \texttt{CALLS}     & Funktion ruft eine andere Funktion auf \\
  \texttt{EXTENDS}    & Klasse erbt von einer anderen Klasse \\
  \texttt{DEFINED\_IN} & Funktion oder Klasse ist in einer bestimmten Datei definiert \\
  \texttt{TESTS}     & Testfunktion testet ein bestimmtes Symbol \\
  \texttt{DOCUMENTS} & Docstring oder Kommentar dokumentiert ein Symbol \\
  \bottomrule
\end{tabularx}

\vspace{0.5em}

Ergänzend werden Git-History-Metriken als Knotenattribute gespeichert: Änderungshäufigkeit und letzte Commit-Nachricht. Diese ermöglichen risikogewichtetes Ranking: Häufig modifizierte Dateien werden bei gleicher Hop-Distanz bevorzugt.

\subsection*{Schicht 2: Living Graph mit inkrementeller Git-Synchronisation}

Ein Git post-commit Hook analysiert nach jedem Commit, welche Dateien sich verändert haben. Nur die betroffenen Knoten und ihre direkten Kanten werden neu geparst und aktualisiert; ein vollständiger Rebuild entfällt. Damit bleibt der Graph dauerhaft synchron mit dem Codestand, ohne dass manuelle Wartung nötig ist. Dies ist die \emph{Living}-Eigenschaft des Systems: Der Wissensgraph ist kein statisches Dokument, sondern ein kontinuierlich aktualisiertes Abbild der Codebase.

\subsection*{Schicht 3: Proaktive Context Assembly}

Für eine eingehende Aufgabenbeschreibung läuft folgende Pipeline:

\begin{enumerate}
  \item \textbf{Entity Extraction:} Aus dem Aufgabentext werden genannte Funktionsnamen, Dateinamen und Schlüsselbegriffe extrahiert (via kleinem LLM-Call oder Regex-Matching).
  \item \textbf{Seed-Node-Identifikation:} Die extrahierten Entitäten werden im Graphen als Startknoten lokalisiert.
  \item \textbf{Graph-Traversal:} Von den Startknoten aus wird der Graph in konfigurierbarer Hop-Tiefe traversiert. Unterschiedliche Traversal-Strategien sind wählbar: Strukturell (Aufruf- und Importkanten), historisch (häufig co-modifizierte Knoten) und testbasiert (zugehörige Testknoten).
  \item \textbf{Ranking und Filterung:} Knoten werden nach Relevanz geordnet (Hop-Distanz, Testabdeckung, Risikogewicht) und auf ein konfigurierbares Tokenbudget begrenzt.
  \item \textbf{Serialisierung:} Der Teilgraph wird als strukturierter Kontext-Block (Markdown oder JSON) ausgegeben, inklusive expliziter Angaben, \emph{warum} jeder Knoten aufgenommen wurde.
\end{enumerate}

\begin{figure}[htbp]
\centering
\begin{tikzpicture}[
  node distance=0.45cm,
  pbox/.style={
    draw, rounded corners=3pt,
    minimum width=2.15cm, minimum height=1.1cm,
    align=center, inner sep=4pt,
    font=\small\bfseries, fill=orange!12
  },
  exlbl/.style={font=\scriptsize, text=gray!70, align=center},
  arr/.style={-{Stealth[length=5pt,width=4pt]}, thick, gray!55}
]

%% Pipeline boxes
\node[pbox] (e1) {Entity\\Extraction};
\node[pbox, right=of e1] (e2) {Seed-\\Lookup};
\node[pbox, right=of e2] (e3) {Graph-\\Traversal};
\node[pbox, right=of e3] (e4) {Ranking\\\& Filter};
\node[pbox, right=of e4] (e5) {Seriali-\\sierung};

%% Arrows between boxes
\draw[arr] (e1) -- (e2);
\draw[arr] (e2) -- (e3);
\draw[arr] (e3) -- (e4);
\draw[arr] (e4) -- (e5);

%% Input / output labels above
\node[exlbl, above=5pt of e1] {\textit{Aufgabe}};
\node[exlbl, above=5pt of e5] {\textit{Kontext-Block}};

%% Example row below
\node[exlbl, below=8pt of e1] {\glqq{}fix cart\\checkout\grqq{}};
\node[exlbl, below=8pt of e2] {\{checkout,\\cart, Cart\}};
\node[exlbl, below=8pt of e3] {32 Knoten,\\2 Hops};
\node[exlbl, below=8pt of e4] {Top 15,\\sortiert};
\node[exlbl, below=8pt of e5] {Markdown\\+~Pfad-Audit};

%% Dotted separator line
\draw[dotted, gray!50]
  ([xshift=-0.3cm, yshift=-0.35cm]e1.south west) --
  ([xshift= 0.3cm, yshift=-0.35cm]e5.south east);

\end{tikzpicture}
\caption{Fünfstufige Context-Assembly-Pipeline (\texttt{get\_context}). Die untere Zeile
zeigt ein Beispiel für die Aufgabe \emph{,,fix cart checkout``}.}
\label{fig:pipeline}
\end{figure}


\subsection*{Transparenz als First-Class-Feature}

Jeder ausgelieferte Kontext-Block enthält einen \emph{Entscheidungspfad}: Eine maschinenlesbare Begründung, über welche Kanten jeder Knoten erreicht wurde (\texttt{checkout() aufgenommen via CALLS-Kante von cart.ts, Zeile 42}). Entwickler können diesen Teilgraphen vor Agentausführung inspizieren und gezielt erweitern oder beschneiden. Nach Abschluss der Aufgabe dient der Entscheidungspfad als Audit-Trail.

\subsection*{Integration via MCP}

Das System wird als \emph{MCP-Server} implementiert und stellt drei Tools bereit:

\begin{description}
  \item[\texttt{get\_context(task)}] Assembliert und gibt den Kontext-Block für eine Aufgabenbeschreibung zurück.
  \item[\texttt{inspect\_graph(node\_id)}] Gibt den direkten Nachbarschaft-Subgraphen eines Knotens zurück.
  \item[\texttt{update\_graph()}] Triggert manuell eine Graph-Aktualisierung (außerhalb des automatischen Git-Hooks).
\end{description}

Claude Code, Cursor und andere MCP-kompatible Clients können diese Tools direkt aufrufen. Damit ist das System ohne Änderung am Agent selbst nutzbar.

\subsection*{Stand der Vorarbeiten}

\begin{tcolorbox}[colback=blue!4!white, colframe=blue!40!black, title=Proof-of-Concept bereits implementiert, fonttitle=\bfseries]
Im Rahmen der Einarbeitungsphase wurde ein funktionsfähiger Proof-of-Concept für die TypeScript-Komponente entwickelt und auf dem eigenen WebshopTemplate-Projekt validiert:

\begin{itemize}[noitemsep]
  \item \textbf{Parser (M2):} tree-sitter-Parser für TypeScript/TSX extrahiert Datei-, Funktions- und Klassenknoten sowie IMPORTS-, CALLS-, EXTENDS- und DEFINED\_IN-Kanten.
  \item \textbf{Graph (M3):} Neo4j-Instanz mit 59 Dateiknoten, 64 Funktionsknoten und 25 aufgelösten CALLS-Kanten; Git-History-Metriken (\texttt{commitCount}, \texttt{lastCommitMessage}) als Knotenattribute; Performance-Indexe.
  \item \textbf{Context Assembly (M5):} Fünfstufige Pipeline (\texttt{get\_context}) mit Entity Extraction, Seed-Node-Identifikation, konfigurierbarer Graphtraversierung, Relevanz-Ranking und Markdown-Serialisierung inklusive Entscheidungspfad.
  \item \textbf{MCP-Server (M6):} Drei Tools (\texttt{get\_context}, \texttt{inspect\_graph}, \texttt{update\_graph}) als MCP-Server implementiert und in Claude Code registriert.
\end{itemize}

Die Living-Eigenschaft (automatischer Git-Hook, M4) und die Python-Komponente sind in der eigentlichen Masterarbeit zu entwickeln.
\end{tcolorbox}
